\scanpage{27}
\mspage{8}
We take an automorphism \inserted{$u$} of the groupoid $\widehat{\overline{\underline{\mathrm{II}}}}_J$ which \inserted{commutes with $G_J$ and which} fixes the vertices, and which as above leaves invariant the subgroupoids $\widehat{\overline{\underline{\mathrm{II}}}}_{J,i}$ $(i\in J)$, \deleted{the groupoids} $\overline{\underline{\mathrm{II}}}_{J,i}$ generated by the \inserted{2} vertices $R_i^\omega$ $(\omega\in\uline{\omega})$ and the two arrows $a_i^\omega$ $(\omega\in\uline{\omega})$.
\deleted{Then} The endomorphism $u$ is known when one knows the $u(a_i^\omega)$, which are of the form
\[
(20)\qquad u(a_i^\omega)=a_i^\omega\varphi_i^\omega\bigl(\underbrace{l_0^\gamma}\bigr)
\qquad\begin{array}{l}\gamma=\frac{p-1}{2}\in\widehat{\mathbb Z}\\p\in\widehat{\mathbb Z}^{*}\end{array}
\]
\margin{NB: $\lambda\in\widehat\Pi_{0,3}$, no need to require at the outset that $\lambda$ of the form $l_0^\gamma\ldots$}
\uline{and} the
\[
(21)\qquad u(c_i^\omega)=\varphi_{\omega'i}^{\omega'}(\eta)c_i^\omega=c_i^\omega\psi_i^\omega(\eta)
\qquad\eta\in\widehat\Pi_{0,3}\quad\text{independent of }i,\omega
\]
We shall then have (compare $(11)$)
\[
(22)\qquad u(b_i^\omega)=b_i^\omega\varphi_{\omega i}^\omega(\beta)\qquad\text{with}
\]
\[
(23)\qquad\beta=\sigma_\infty(\eta)\eta^{-1}
\]
\margin{NB becomes prettier if $\eta$ is replaced by its $\sigma_\infty\eta$, or $\beta$ by $\beta^{-1}$}
The relations to be satisfied \deleted{\unsure{are}} then already written in \inserted{$p$ (or $\gamma$) and} $\beta$, \deleted{\unsure{the}} in $(12)$ --- but the first relation $\beta\sigma_\infty(\beta)=1$ a consequence of $(23)$, there remains the second which is therefore written
\margin{\uline{NB} If $u$ replaced by $u'=\tau(u)$, then $\eta$ is replaced by $\eta'=\tau_\infty(\eta)$, $\beta$ by $\tau_\infty(\beta)$, but $\lambda$ by $\varepsilon_0^{-1}(\tau_\infty(\lambda^{-1}))$.}
\[
(24)\qquad\deleted{(l_1^{-\gamma}\sigma_\infty(\eta)\eta^{-1}}
\quad\alpha\rho^{-1}(\alpha)\rho^{-2}(\alpha)=1
\qquad\text{with}\quad\alpha=\underbrace{l_1^\gamma}_{\sigma_\infty(\lambda)}\underbrace{\sigma_\infty(\eta)\eta^{-1}}_\beta.
\]
We shall have here
\[
(25)\qquad
u(\widetilde a_i^\omega)=\psi_{\omega\deleted{\unsure{'}}i}^{\omega\deleted{\unsure{'}}}\bigl(\underbrace{\eta^{-1}l_0^\gamma\varepsilon_0(\eta)}_{\widetilde\alpha=\eta^{-1}\varepsilon_0(\lambda\eta)}\bigr)\widetilde a_i^\omega
=\widetilde a_i^\omega\psi_{\omega'i}^{\omega'}\bigl(\underbrace{\varepsilon_0^{-1}(\eta^{-1})l_0^\gamma\eta}_{\varepsilon_0^{-1}(\widetilde\alpha)}\bigr)
\]
\inserted{valid for any $\lambda$}
\margin{$\widetilde\alpha=\eta^{-1}l_0^\gamma\varepsilon_0(\eta)$}
and
\[
(25\text{ bis})\quad\deleted{\unsure{\alpha=}}\quad\alpha=\rho\bigl(\eta\widetilde\alpha\rho^{-1}(\eta)^{-1}\bigr)
\]
\margin{valid for arbitrary $\lambda$, not necessarily of the form $l_0^\gamma$.}
\deleted{for} in the equation $(24)$ in $\alpha$ (invariant under substitutions $\alpha\mapsto\rho^{-1}(\alpha)$ (or $\alpha\mapsto\rho\alpha=\rho^{-2}(\alpha)$) and $\alpha\mapsto x\alpha\rho^{-1}(x)^{-1}$,
\margin{26}
